% Originally: content/week-13.tex, \section{Submanifolds with boundary and induced orientation} (Corsin Week 13)


\section{Submanifolds with boundary and induced orientation}
\label{sec:submanifolds_with_boundary}

% Source: Corsin Nick/Class Notes/Week 13.pdf, pp. 6--7
\begin{ainote}
Corsin recaps \cref{thm:local_parametrization} verbatim --- for every $p \in M$, an
open neighbourhood $V$ of $p$, an open set $U \subseteq \mathbb{R}^k$, and $\phi : U \to
V\cap M$ that is smooth, has $D\phi_x$ injective for all $x \in U$, and is a homeomorphism onto
$V\cap M$ --- before extending it to submanifolds with boundary below. Not repeated; see
\cref{sec:submanifolds}.
\end{ainote}

\begin{definition}[Submanifold with boundary]
\label{def:submanifold_with_boundary}
For a $k$-dimensional submanifold \emph{with boundary}, we require similarly that for a local
parametrization $\phi : U \to V\cap M$ as in \cref{thm:local_parametrization}, either $U
\subseteq \mathbb{R}^k$ is open, or
\[U = U' \cap \mathcal{H}^k, \qquad U' \subseteq \mathbb{R}^k \text{ open}, \qquad
  \mathcal{H}^k := \{(x_1,\ldots,x_k) \in \mathbb{R}^k : x_1 \leq 0\}\]
the (closed) \newterm{left half-space}. If $p \in \partial M$ and $\phi : U \subseteq
\mathcal{H}^k \to V\cap M$ is a local parametrization of $M$ around $p$, then $\phi|_{\partial
\mathcal{H}^k}$ (identifying $\partial\mathcal{H}^k = \{0\}\times\mathbb{R}^{k-1} \cong
\mathbb{R}^{k-1}$) is a local parametrization of $\partial M$ around $p$.
\end{definition}

% Sonnet 5 (Medium) --- FIG-W13-01
\begin{center}
\begin{tikzpicture}[>=Latex, scale=1.0]
  \draw[green!50!black, thick] (0,0) .. controls (0.3,1.3) and (1.3,2.0) .. (2.6,1.9)
    .. controls (3.9,1.8) and (4.6,1.0) .. (4.4,0.1) .. controls (3.5,-0.5) and (1.5,-0.6)
    .. (0,0);
  \fill (0.9,0.55) circle (1.1pt);
  \draw[red, dashed, thick] (-0.05,-0.55) .. controls (0.4,0.0) and (0.5,1.1) .. (0.15,1.55);
  \draw[red, thick, ->] (0.9,0.55) -- (1.75,0.75) node[right, font=\scriptsize] {$\partial_2\phi$};
  \draw[red, thick, ->] (0.9,0.55) -- (0.55,1.3) node[above, font=\scriptsize] {$\partial_1\phi$};
  \node[font=\scriptsize] at (0.6,-0.85) {$p \in \partial M$};
  \draw[->] (5.3,0.7) -- (6.3,0.7) node[right, font=\scriptsize] {};
  \node[font=\scriptsize] at (5.8,1.0) {$\phi$};
  \draw[red, thick] (7.1,0.1) rectangle (8.4,1.4);
  \draw[red, thick, dashed] (7.1,0.1) -- (7.1,1.4);
  \draw[red, ->] (7.1,0.75) -- (7.7,0.75) node[right, font=\scriptsize] {$e_1$};
  \draw[red, ->] (7.1,0.75) -- (7.1,1.25) node[above, font=\scriptsize] {$e_2$};
  \node[font=\scriptsize] at (7.1,-0.15) {$0$};
  \node[font=\scriptsize] at (7.75,-0.2) {$U \subseteq \mathcal{H}^2$};
\end{tikzpicture}
\end{center}

\begin{definition}[Induced orientation]
\label{def:induced_orientation}
If $\{\partial_1\phi,\ldots,\partial_k\phi\}$ is the oriented basis of $T_pM$ given by a local
parametrization $\phi$ around $p \in \partial M$, then $\{\partial_2\phi,\ldots,\partial_k
\phi\}$ is the \newterm{induced} oriented basis of $T_p(\partial M)$ around $p$.
\end{definition}

\subsection{Example: cylinder}
\label{sec:cylinder_orientation_example}

% Source: Corsin Nick/Class Notes/Week 13.pdf, pp. 8--9
Consider the cylindrical surface
\[C = \{(x,y,z) \in \mathbb{R}^3 : y^2+z^2 = r^2,\ 0\leq x\leq L\},\]
a $2$-dimensional submanifold with boundary, and the two parametrizations $\phi,\psi :
[-L,0]\times[0,2\pi] \to C$ (both domains $\subseteq \mathcal{H}^2$),
\[\phi(x,\theta) = \begin{pmatrix} -x \\ r\cos\theta \\ -r\sin\theta \end{pmatrix}, \qquad
  \psi(x,\theta) = \begin{pmatrix} x+L \\ r\cos\theta \\ r\sin\theta \end{pmatrix}.\]
Both cover all of $C$: as the local $x$-coordinate ranges over $[-L,0]$, $-x$ and $x+L$ both
range over $[0,L]$.

The two are positively oriented with respect to each other: writing out the composition,
\[\phi^{-1}\circ\psi(x,\theta) = \phi^{-1}(x+L,\,r\cos\theta,\,r\sin\theta) = (-x-L,-\theta)
  = -(x+L,\theta),\]
so $D(\phi^{-1}\circ\psi) = \begin{pmatrix}-1&0\\0&-1\end{pmatrix}$ and $\det D(\phi^{-1}
\circ\psi) = 1 > 0$: $\{\phi,\psi\}$ is a valid orientation of $C$ by
\cref{def:orientable_manifold}.

But restricted to the line $\partial\mathcal{H}^2 = \{(x,\theta) : x=0\}$, $\phi$ and $\psi$
parametrize \emph{opposite} boundary components of the cylinder --- $\phi(0,\cdot)$ traces the
circle at $x=0$, $\psi(0,\cdot)$ the circle at $x=L$ --- and \cref{def:induced_orientation}
gives them opposite rotational senses when both are drawn in the same $(y,z)$-view:

% Sonnet 5 (Medium) --- FIG-W13-02
\begin{center}
\begin{tikzpicture}[>=Latex, scale=0.95]
  \draw[blue!70!black, thick] (0,2.4) rectangle (4.5,3.5);
  \draw[blue!70!black, thick] (0,2.4) circle (0.55 and 0.55);
  \draw[blue!70!black, thick] (4.5,2.95) ellipse (0.35 and 0.55);
  \draw[red, thick, ->] (-0.5,3.5) arc (150:490:0.55);
  \node[red, font=\scriptsize] at (-0.85,3.5) {$t\mapsto\phi(0,t)$};
  \node[font=\scriptsize] at (0,1.7) {\textbf{left end} ($x=0$)};

  \draw[blue!70!black, thick] (6.5,2.4) rectangle (11,3.5);
  \draw[blue!70!black, thick, dashed] (6.5,2.95) ellipse (0.35 and 0.55);
  \draw[blue!70!black, thick] (11,2.4) circle (0.55 and 0.55);
  \draw[orange, thick, ->] (11.5,2.4) arc (-30:330:0.55);
  \node[orange, font=\scriptsize] at (11.9,3.5) {$t\mapsto\psi(0,t)$};
  \node[font=\scriptsize] at (8.75,1.7) {\textbf{right end} ($x=L$)};
\end{tikzpicture}
\end{center}

\begin{ainote}
This is not a contradiction, and Corsin does not present it as one: $\{\phi,\psi\}$ is a
perfectly good orientation of $C$ as a whole (the overlap check above is what
\cref{def:orientable_manifold} demands, and it passes). What the picture shows is that a
\emph{consistently} oriented cylinder induces boundary orientations on its two ends that look
like opposite rotational senses when both circles are drawn from the same external viewpoint
--- exactly the familiar fact, from the divergence theorem, that the outward normal at the two
ends of a solid cylinder points in opposite $x$-directions. Induced orientation is doing its
job correctly here, not failing.
\end{ainote}
